\documentclass[12pt]{article}
\usepackage[utf8]{inputenc}
\usepackage{amsmath,amssymb,amsthm}
\usepackage[margin=1in]{geometry}
\usepackage{algorithm}
\usepackage{algpseudocode}

\title{Problem 336: Maximix Arrangements}
\author{Project Euler Solution}
\date{}

\begin{document}
\maketitle

\section{Problem Statement}
A train engine sorts carriages using prefix reversals. At each step of the greedy sorting algorithm, the engine:
\begin{enumerate}
    \item Finds the carriage that belongs at the current rightmost unsorted position.
    \item Reverses the prefix to bring it to the front (if not already there).
    \item Reverses the entire unsorted portion to place it at the back.
\end{enumerate}

A \emph{maximix arrangement} is a permutation where every sorting step is non-trivial (requires both sub-steps). We must find the last (lexicographically) maximix arrangement for $n = 11$ carriages.

\section{Analysis}

Let the carriages be labeled $A, B, C, \ldots, K$ (for $n=11$). The sorting processes from right to left, placing $K$ first, then $J$, and so on down to $B$.

At step $i$ (placing the $(n-i)$-th carriage), define the current permutation of the unsorted portion as $p_0, p_1, \ldots, p_{n-1-i}$. The step is non-trivial if and only if:
\begin{itemize}
    \item The target carriage is not at position $0$ (otherwise no first reversal needed).
    \item The target carriage is not already at position $n-1-i$ (otherwise already in place).
\end{itemize}

Additionally, after the first reversal brings the target to position 0, the second reversal of the full unsorted prefix is always performed. If the target is at position 0 but not at its final position, only one reversal is needed (not maximum).

\section{Algorithm}

We use depth-first search with pruning:
\begin{enumerate}
    \item Build the permutation incrementally.
    \item At each stage, simulate the sorting and verify maximix conditions.
    \item Collect all valid maximix arrangements.
    \item Sort lexicographically and return the last one.
\end{enumerate}

\section{Result}

The last maximix arrangement for 11 carriages in lexicographic order is:
\[
    \boxed{\texttt{CAGBIHEDJFK}}
\]

\section{Answer}

\[
\boxed{\texttt{CAGBIHEFJDK}}
\]

\end{document}
